\fsection{Posibles ampliaciones}
El alcance del proyecto es la celda de bin picking descrita. Sobre ella se plantean las siguientes ampliaciones,
que quedan como trabajo futuro:

\begin{itemize}
	\item Identificación y clasificación de piezas: trabajar con piezas de distintos tipos en el mismo contenedor
	      y depositar cada una en una bandeja según su tipo, usando las bandejas 2 y 3 que ya están previstas en la
	      mesa. Para ello se añade una estación de lectura fija junto a la celda: tras el agarre, el robot presenta
	      la pieza ante el lector, el PLC recibe su identificador y decide en qué bandeja la deposita. La identificación puede hacerse con un lector de códigos QR o DataMatrix con
	      Profinet (por ejemplo, SIMATIC MV500), con etiquetas en la base o en un lateral de la pieza, o con un lector
	      RFID como el SIMATIC RF260R, que no necesita visión directa del tag. En ambos casos basta con un FB nuevo en
	      el programa del PLC y un punto de paso más en el programa del robot, sin modificar el sistema de Siemens.
	\item Eje traslacional: añadir un séptimo eje lineal al robot para ampliar su alcance y trabajar con varios
	      contenedores.
	\item Reentrenamiento del modelo de IA con piezas propias del entorno productivo, para mejorar la tasa de
	      detección con geometrías distintas a las del modelo preentrenado.
\end{itemize}
